\documentclass[12pt]{article}
\usepackage{amsmath, amssymb}
\usepackage{geometry}% 1-inch margins
\usepackage{listings}
\geometry{letterpaper, margin=1in}% 1-inch margins

\title{MATH 374 (Discrete Structures)\\Chapter 1.5 Homework}
\author{Jacob Vaught}

\begin{document}
\maketitle

\section*{EXAMPLE 39}

Suppose that a Prolog database contains the following entries:
\newline
\textit{eat}(bear, fish).\newline
\textit{eat}(fish, little-fish).\newline
\textit{eat}(little-fish, algae).\newline
\textit{eat}(raccoon, fish).\newline
\textit{eat}(bear, raccoon).\newline
\textit{eat}(bear, fox).\newline
\textit{eat}(fox, rabbit).\newline
\textit{eat}(rabbit, grass).\newline
\textit{eat}(bear, deer).\newline
\textit{eat}(deer, grass).\newline
\textit{eat}(wildcat, deer).\newline
\textit{animal}(bear).\newline
\textit{animal}(fish).\newline
\textit{animal}(little-fish).\newline
\textit{animal}(raccoon).\newline
\textit{animal}(fox).\newline
\textit{animal}(rabbit).\newline
\textit{animal}(deer).\newline
\textit{animal}(wildcat).\newline
\textit{plant}(grass).\newline
\textit{plant}(algae).\newline
\textit{prey}(x) if \textit{eat}(y, x) and \textit{animal}(x).\newline

\newpage
\subsection*{Exercise 1.5.7}

Formulate a Prolog rule that defines “herbivore” to add to the database of Example 39.

\textbf{Solution:}

\textit{herbivore}(y) if \textit{animal}(y) and not(\textit{eat}(y, z) and \textit{animal}(z)) and \textit{eat}(y, x) and \textit{plant}(x).

\subsection*{Exercise 1.5.8}

If the rule of Exercise 1.5.7 is included in the database of Example 39, what is the response to the following query?

\textbf{Query:}
Which(x: \textit{herbivore}(x))?


\textbf{Solution:}
Based on the rule defined in Exercise 1.5.7, the animals that are herbivores are those that do not eat other animals and do eat plants. From the database, we can conclude that the only animals that meet these criteria are the rabbit and the deer. Therefore, the response to the query will be:
\begin{lstlisting}
rabbit
deer
\end{lstlisting}
\end{document}
